% Section 4 — Evaluation methodology
% Page budget: ~2 pages in acmart sigconf two-column (~1700 words).
% Source-of-truth: DOCS/docs8/PAPER-FINDINGS.md, RQ-CLOSURE-TABLE.md,
% per-section SUMMARY.md, and DOCS/docs7/REAL-DATA-WoL-PLAN.md.
%
% Editorial discipline:
%   - No measured numbers here. Those live in §5--§6. Only definitions,
%     splits, baselines, and protocol.
%   - All three datasets are described in this section; §3 referred
%     forward here. The construction-details paragraphs are
%     intentionally concrete (record counts, file structure) because
%     reviewers asking ``what is the dataset?'' should not have to
%     hunt across the paper.

\section{Evaluation Methodology}
\label{sec:evaluation}

This section fixes the methodology under which every measurement in \S\ref{sec:results}--\S\ref{sec:cross-app} is produced. We describe the three datasets (\S\ref{sec:evaluation:datasets}), the train/validation/test partitions and the visibility rule that prevents test-time leakage (\S\ref{sec:evaluation:splits}), the metrics we report (\S\ref{sec:evaluation:metrics}), the baselines the agent is compared against (\S\ref{sec:evaluation:baselines}), the statistical envelope (\S\ref{sec:evaluation:stats}), and the implementation details a reviewer needs to reproduce a result (\S\ref{sec:evaluation:impl}).

\subsection{Datasets}
\label{sec:evaluation:datasets}

The agent is evaluated on three datasets chosen to stress different parts of its design: a controlled synthetic microservices benchmark, a structurally different polyglot synthetic application, and a corpus of real Apache Jira tickets paired with telemetry-style queries.

\subsubsection{Synthetic microservices telemetry (\dsob{})}
\label{sec:evaluation:datasets:ob}

The first corpus extends Google's Online Boutique~\cite{onlineboutique}, a ten-service e-commerce demo whose services are written in six languages and communicate over remote procedure calls with a cache backing the cart service. We forked the upstream code and added production-realistic telemetry under a single binding rule: \emph{anything a real production service would not emit must not be added}. We reject any field that names our scenario taxonomy---no scenario identifiers in span attributes, no fault-type strings in log fields, no alert names tied to our incident families. The discipline makes debugging harder but the resulting telemetry indistinguishable from what an instrumented production fleet would emit; the evaluation that follows describes the production task rather than a fictional easier surrogate.

Three telemetry channels are added to every service. The logging channel emits a structured record per remote-procedure call, plus a structured error record at every dependency boundary capturing what was being called, with what operation, with what error class, and at what retry attempt. A small set of business events records domain-level signal (cart-size changed, order placed, payment charged) without ever encoding scenario identity. The tracing channel attaches errors to the span where they occurred and adds manual child spans for dependency calls with standard semantic attributes~\cite{otel-spec}. The metrics channel exports the rate, error rate, and latency of every handler, together with per-dependency client breakdowns and standard runtime gauges from each language's default exporter.

Faults are injected through four conceptual mechanisms that collectively define a taxonomy of twenty-seven scenario families. \emph{Service removal} takes a chosen service offline, backing the eight hard-outage families and the cart-dependency outage. \emph{Behavioural perturbation} changes the configuration of a still-running service (slowing it, biasing its responses, raising its traffic rate), backing latency-injection and traffic-pressure scenarios. \emph{Pod-level disruption} terminates and replaces pods to model graceful redeploys, single-pod restarts within a healthy replica set, and pathological flapping. \emph{Network-level disruption} delays, drops, or partitions traffic between specific service pairs via a chaos-engineering tool~\cite{chaosmesh}. A fifth, trivial mechanism is passive observation under normal traffic, which produces the negative-class baseline windows.

Each fault produces an \emph{incident episode} with three associated \emph{telemetry windows}: a pre-fault baseline captured just before injection, an active-fault window during the injection itself, and a recovery window after the fault is cleared. Long-running scenarios such as slow memory leaks add a fourth, longer observation window. The full corpus contains $6{,}720$ labeled telemetry windows from $100$ independent fault-injection runs, of which $1{,}008$ are held out for the test split (\S\ref{sec:evaluation:splits}). Each window carries five categories of label that downstream training and evaluation consume: a triage class (ticket-worthy, borderline, noise); for ticket-worthy windows, the Jira metadata the corresponding ticket would carry (severity, affected components, reason class); a memory-match ground-truth list of past tickets whose scenario family matches the window's, restricted to tickets time-visible at the window's timestamp and originating from a different fault-injection run; a Boolean novelty flag that fires when the match list is empty; and a fixed-dimensional vector of pre-computed numeric features summarising the window's logs, traces, metrics, and Kubernetes-state counters. A separate $347$-ticket Jira memory corpus, humanised from the scenario taxonomy into engineer-voice prose grounded in TAWOS length and code-block distributions, is paired with the telemetry; a $110$-ticket distractor pool is used to evaluate retrieval robustness under memory contamination.

\subsubsection{OpenTelemetry Demo (\dsotel{})}
\label{sec:evaluation:datasets:otel}

The second corpus is built from the OpenTelemetry Demo~\cite{otel-demo}, also known as the Astronomy Shop: a structurally different polyglot e-commerce application with fifteen-plus services across twelve-plus languages, Kafka-based asynchronous communication, and a feature-flag service that lets us inject many of the same conceptual faults without writing per-service patches. The lab uses the same instrumentation discipline as \dsob{}, the same conceptual fault mechanisms, and the same window structure, but inhabits a different application. The result is $1{,}643$ telemetry windows from twenty-two fault-injection runs across fifty-three scenario families (a broader taxonomy than \dsob{}'s twenty-seven), with $247$ held out for the test split. The paired Jira corpus is smaller---$147$ humanised tickets---and the per-family memory depth is correspondingly thinner, which stresses retrieval primitives that depend on memory density. The \dsotel{} corpus is the controlled-but-different cross-application test for the agent's transfer claim: same conceptual setup, different application underneath.

\subsubsection{Real Apache Jira tickets (\dswol{})}
\label{sec:evaluation:datasets:wol}

The third corpus is the strongest external-validity test. It draws on $2{,}000$ real Jira tickets from seven Apache Software Foundation open-source projects: Kafka, MariaDB Server, Cassandra, Spark, ZooKeeper, HDFS, and Solr. Tickets are selected from each project's public bug tracker, filtered for valid resolution dates and non-trivial issue bodies, and aligned to a telemetry-style query format by treating the ticket's symptom paragraph (when present) as the synthetic ``window evidence text'' that a triage system would see at incident time, with the rest of the ticket (diagnosis, root-cause discussion, fix) reserved as the memory-side text the retriever's gold labels point to. The dataset's gold-match relation is defined at two levels: a \emph{coarse} relation (same project family) and a \emph{strong} relation (same project family plus symptom-overlap above a configured threshold).

The test split holds out two entire project families (Kafka and MariaDB Server) so that no test window's gold ticket comes from a project the model has seen during fine-tune. This is a family-held-out evaluation, the strictest cross-domain setting~\cite{tawos}; it tests whether the agent's retrieval ranking carries semantic structure that generalises beyond the training projects' vocabulary or whether it has merely memorised them. The split has $304$ test windows, of which $55$ are evaluable for retrieval (have at least one visible gold ticket under the time-ordered visibility rule). The smaller evaluable subset reflects the dataset's heavier holdout: by construction most test windows belong to projects whose ticket history is not visible at retrieval time. Confidence intervals on \dswol{} numbers are correspondingly wider; we report this honestly in \S\ref{sec:cross-app}.

\subsection{Splits and the visibility rule}
\label{sec:evaluation:splits}

For each dataset, windows are partitioned chronologically and stratified along the scenario family (or project family, for \dswol{}) into a training set on which all models are fit, a validation set used to tune thresholds and select operating points, and a held-out test set on which every headline number in this paper is computed. On the synthetic datasets every scenario family appears in all three sets, but each set draws windows from disjoint fault-injection runs of that family; this matches a mature on-call team's operational view, in which past tickets exist for most failure modes the team encounters but each present-day incident still represents a previously-unseen instance of its family. On the real-Jira dataset the split is stricter, holding out two project families entirely from training to test cross-project transfer.

A complementary partition on \dsob{} holds out an entire scenario family for the out-of-distribution analysis. Twenty-seven such folds are constructed, one per family; each fold trains on the remaining twenty-six and evaluates on the held-out family alone. This is a strictly harder regime than the in-distribution split because the model has never seen any examples of the held-out family during fit, and the resulting per-family novelty score measures whether the system has learned a generic notion of \emph{``no past ticket matches''} or has merely memorised the families it saw at training time.

A strict time-ordered visibility rule is enforced across both partitions. A test window with timestamp $t$ may retrieve only memory tickets whose own creation timestamp is strictly less than $t$, and is additionally prevented from seeing any ticket that originated in the same fault-injection run as the window itself. The first rule prevents cheating by ``predicting'' future incidents; the second prevents the subtler form of leakage in which a window retrieves a ticket filed about its own incident. Both rules are enforced at retrieval time by the corpus loader, not optionally and not at training time.

\subsection{Metrics}
\label{sec:evaluation:metrics}

We report five families of metric. Each family answers a different question.

\paragraphhead{Retrieval.} The primary retrieval metric is Hit@5: among the test windows for which the dataset records at least one gold past ticket, the fraction whose top-five ranked suggestions contain a gold. Hit@1 is the same with a top-one cutoff. Mean reciprocal rank averages the inverse rank of the first gold appearing in the system's output (zero when none of the top results is gold), rewarding systems that rank the right ticket high rather than merely surfacing it somewhere. Hit@10 is reported for context.

\paragraphhead{Triage detection.} Triage accuracy is the fraction of test windows on which the system's emitted triage decision matches the gold triage label (ticket-worthy versus noise after collapsing the borderline class). We additionally report the area under the precision-recall curve, the operating-point precision at a $1\%$ false-positive-rate budget, and the expected calibration error.

\paragraphhead{Novelty.} Of the windows the system flagged as novel, what fraction truly had no past match (novel precision)? Of the truly-novel windows, what fraction did the system correctly flag (novel recall)? The harmonic mean is the novelty F1.

\paragraphhead{Cost.} We report four cost variables: total wall time per test split, total language-model tokens emitted, total dollar-equivalent cost (computed under a per-skill cost table; see Appendix~A in PAPER-FINDINGS.md and \S\ref{sec:results:cost} for the assumptions), and the percentage saved relative to an always-on cascade counterfactual in which every diagnostic primitive runs on every window. The savings percentage is invariant to the underlying language-model unit cost, because the counterfactual and the actual scale together with that unit cost.

\paragraphhead{Operations.} Pages-per-incident is the number of paged tickets divided by the number of distinct incident identifiers in the test split. Suppressions-fired is the count of state-suppression events emitted during the run (a positive integer means the cross-window state layer was exercised, not bypassed). Distinct plan identifiers is the number of structurally distinct plans the controller emitted across the test split, with a fixed-pipeline system scoring one and an adaptive one scoring higher; the metric measures whether the controller's branching machinery is actually exercised on real inputs.

\paragraphhead{Failure-mode catalog.} For the on-demand evidence-gathering loop, we report the per-tool counts of the five failure modes from \S\ref{sec:approach:tooluse}: success, empty result, looping repeated call, budget exhausted, tool error. A reviewer asking ``what about robustness?'' should see this catalog rather than an aggregated success rate.

\subsection{Baselines}
\label{sec:evaluation:baselines}

The agent is compared against three baselines.

\paragraphhead{Sparse-only retrieval.} A BM25~\cite{bm25} retriever indexed on the same memory corpus. BM25 is the naive baseline a reviewer expects: ``can your system beat keyword matching?'' We report Hit@K for BM25 on the agent-evaluable subset of each test split so that the comparison is apples-to-apples (same windows, same gold relation).

\paragraphhead{Always-on cascade counterfactual.} A virtual baseline in which every diagnostic primitive runs on every window, regardless of the controller's gates. The cost is computed by summing the per-primitive cost (from the per-skill cost table) across every window in the test split, including primitives the agent's controller would have skipped. The retrieval and triage quality of this counterfactual is identical to the agent's at the active-fault branch's full skill set, because the agent already invokes all of those primitives on uncertain windows. The interesting axis is cost: the gap between the agent's actual cost and the counterfactual is the cost-savings metric.

\paragraphhead{Single-primitive baselines.} For ablation and per-skill attribution we also report each diagnostic primitive's standalone Hit@K when run on its own. The dense bi-encoder retriever, the hybrid sparse-plus-dense retrievers, the structural retriever, and the log-sequence retriever each produce a per-window ranked list that we score independently.

\subsection{Statistical envelope}
\label{sec:evaluation:stats}

Every headline number is reported with a $95\%$ percentile confidence interval obtained from a $1{,}000$-resample paired bootstrap at seed $42$. For pairwise comparisons (agent vs.\ a baseline; ablation cell vs.\ baseline cell) we additionally report a \emph{paired-delta} confidence interval: rather than resampling the test set independently for each system, we resample window identifiers and compute the per-resample delta on the same identifiers in both arms. Paired-delta confidence intervals are substantially tighter than single-report intervals when the two systems agree on most windows---which is the typical pattern for our ablations---and substantially fairer when they do not, because per-window correlations are captured.

The choice of seed and the choice of $1{,}000$ resamples are deliberate. Both are standard for top-tier publication; both are easy to reproduce; and both are the values reviewers expect to see~\cite{rag-lewis}. All bootstrap-CI computations are produced by a single script with the random-state argument exposed, so the entire statistical envelope can be regenerated from the per-window prediction records in the supplementary material.

\subsection{Implementation details}
\label{sec:evaluation:impl}

The agent is implemented as a Python library; the diagnostic primitives are wrapped behind a uniform skill interface. The fine-tuned dense bi-encoder uses sentence-transformers' MiniLM-L6-v2 backbone~\cite{sentencetransformers} with five training epochs at a batch size of sixty-four on a single GPU; fine-tune time is approximately two hours per dataset on a consumer-grade card and reduces to fifteen minutes on the smaller \dsotel{}. The language-model verifier is configured to run against a self-hosted OpenAI-compatible endpoint by default (LM Studio with Qwen-3.6-35B at the time of writing) and against vendor-hosted endpoints (OpenAI, Anthropic, Ollama, vLLM) by changing a single environment variable; all dollar-equivalent cost numbers in this paper are computed at gpt-4o-mini token rates so that the savings percentage is interpretable on a vendor-served deployment. Cache, trace, and per-skill prediction outputs are persisted under a deterministic directory layout that lets a reviewer reproduce any specific window's decision by re-running a single command. The fine-tune random seed, the bootstrap seed, and the controller's gate-threshold defaults are all exposed via command-line arguments.
